\input{../tex-inputs/latex-leseansicht-vorspann}

\section[Gustav Schwarzkopf an Arthur Schnitzler, 2. 6. 1913]{L04257 Gustav Schwarzkopf an Arthur Schnitzler, 2. 6. 1913}
\nopagebreak\mylabel{L04257v}
\rehead{ }\normalsize\beginnumbering\briefempfaengerindex{Schnitzler, Arthur@\textsc{Schnitzler, Arthur}!zzzSchwarzkopf, Gustav@\emph{von Gustav Schwarzkopf}!1913-06-021@{2. 6. 1913}|(be}
\toendnotes[C]{\smallbreak\pagebreak[2]}
\correspDesc{Versand  durch Gustav Schwarzkopf am 2. 6. 1913 in Baden bei Wien
\newline{}Erhalt  durch Arthur Schnitzler im Zeitraum [3. 6. 1913 – 7. 6. 1913?] in Wien}\toendnotes[C]{\smallbreak}
\Standort{CUL, Schnitzler, B 96a.}
\physDesc{Bildpostkarte, 534 Zeichen
\newline{}Handschrift: schwarze Tinte, deutsche Kurrent
\newline{}Versand: Stempel: »\nobreak{}\oindex{Baden bei Wien@\textbf{Baden bei Wien}, \emph{Hauptstadt}|pwk}Baden \textcolor{gray}{2} Niederösterreich, 3. VI. 13, 6\nobreak{}«.  
\newline{}Schnitzler: mit Bleistift Vermerk: »\textsc{Schwarzkop}{[}f{]}« }\toendnotes[C]{\smallbreak}\pstart{}{\pb}Herrn\pend{}\pstart{}D\textsuperscript{r} Arthur Schnitzler\pend{}\pstart{}\textsc{Wien XVIII\oindex{XVIII., Währing@\textbf{XVIII., Währing}, \emph{Verwaltungsgebiet}|pw}}\pend{}\pstart{}Sternwarteſtraße 71\oindex{Wien@\textbf{Wien}!XVIII., Währing@\textbf{XVIII., Währing}!Sternwartestraße 71@\textbf{Sternwartestraße 71}, \emph{Wohngebäude}|pw}.\pend{}{\bigskip}
\pstart
           {\pb}\textcolor{gray}{\textbf{Pension Quisisana}}\oindex{Pension Villa Quisisana@\textbf{Pension Villa Quisisana}, \emph{Beherbergungsgebäude}|pw}\hfill \textcolor{gray}{\textbf{Baden}}\oindex{Baden bei Wien@\textbf{Baden bei Wien}, \emph{Hauptstadt}|pw}\pend
           
\pstart
           \textcolor{gray}{\textbf{Elisabethstrasse 63, 65}}\oindex{Elisabethstraße [Baden bei Wien]@\textbf{Elisabethstraße [Baden bei Wien]}, \emph{Straße}|pw}\hfill \textcolor{gray}{\textbf{N.-Oe.}}\oindex{Niederösterreich@\textbf{Niederösterreich}, \emph{Land}|pw}\pend
           \vspace{1em}
\pstart
           \raggedleft{}{\pb}2. VI 13\pend
           \vspace{0.5em}
\pstart
           Lieber Arthur, ich habe die Abſicht am \uline{5.} Abends in Wien\oindex{Wien@\textbf{Wien}, \emph{Verwaltungsgebiet}|pw} zu{ }ſein. Eigentlich
               wollte ich ſchon morgen fort; aber mein Bruder\pwindex{Schwarzkopf, Max 12.\,6.\,1857 Wien – 14.\,4.\,1928 ebd.@\textsc{Schwarzkopf, Max} (12.\,6.\,1857 Wien – 14.\,4.\,1928 ebd.), \emph{Rechtsanwalt}|pwv} hat mir ſo zugeredet noch zu bleiben. Es iſt auch wirklich{ }ſehr
               hübſch hier. – Es geht mir ganz leidlich, richtiger geſagt, ich habe eigenlich über
               nichts zu klagen, aber es iſt doch noch nicht das Rechte. We{\geminationn} alles gut geht, will ich Sie Freitag
               Nachmittag aufſuchen. Hoffentlich iſt Frau \textsc{Liesel\pwindex{Steinrück, Elisabeth 19.\,11.\,1885 – 7.\,4.\,1920 Partenkirchen@\textsc{Steinrück, Elisabeth} (19.\,11.\,1885 – 7.\,4.\,1920 Partenkirchen)|pw}} da{\geminationn} noch da.\pend
           
\pstart
           Ihr,{\\[\baselineskip]}Ihrer Frau\pwindex{Schnitzler, Olga 17.\,1.\,1882 Wien – 13.\,1.\,1970 Lugano@\textsc{Schnitzler, Olga} (17.\,1.\,1882 Wien – 13.\,1.\,1970 Lugano), \emph{Schauspielerin, Sängerin}|pwv} und Ihnen die herzlichſten{\\[\baselineskip]}Grüße.{\\[\baselineskip]}\spacefill\mbox{Gustav.}\pend
           \leftskip=0em{}\selectlanguage{ngerman}\endnumbering\briefempfaengerindex{Schnitzler, Arthur@\textsc{Schnitzler, Arthur}!zzzSchwarzkopf, Gustav@\emph{von Gustav Schwarzkopf}!1913-06-021@{2. 6. 1913}|)be}\mylabel{L04257h}
\begin{anhang}
\end{anhang}\newcommand{\dateiname}{L04257}\newcommand{\titel}{Gustav Schwarzkopf an Arthur Schnitzler, 2. 6. 1913}\newcommand{\editorInnen}{Herausgegeben von Jahnke, SelmaMüller, Martin Anton}\input{../tex-inputs/latex-leseansicht-abspann}
